\documentclass[12pt]{article}
\usepackage{amsmath, amssymb}
\usepackage[margin=1in]{geometry}
\setlength{\parskip}{6pt}
\title{QUBO Formulation: Hamiltonian Cycle Problem}
\date{}
\begin{document}
\maketitle

\subsection*{Hamiltonian Cycle Problem}

\paragraph{Problem Statement.}
Given an undirected graph $G = (V, E)$ with $N = |V|$ vertices, the Hamiltonian Cycle Problem seeks a cycle that visits every vertex exactly once and returns to the starting vertex. Equivalently, we seek a subset of edges $E' \subseteq E$ such that every vertex has degree exactly two in the subgraph $(V, E')$, and the subgraph forms a single connected cycle rather than disjoint subtours. The objective is to minimize the number of missing cycle edges, where a value of zero indicates the existence of a valid Hamiltonian cycle.

\paragraph{Decision Variables.}
For each pair of distinct vertices $i, j \in V$ with $i < j$, we introduce a binary decision variable $x_{ij} \in \{0, 1\}$. The variable $x_{ij}$ equals $1$ if the edge $(i, j)$ is selected to be part of the cycle, and $0$ otherwise. Since the graph is undirected, variables are defined only for $i < j$ to avoid duplication.

\paragraph{QUBO Derivation.}
Step 1 --- Write the original objective in general notation. The objective minimizes the number of missing edges. For a graph with $N$ vertices, a valid cycle contains exactly $N$ edges. Thus, the objective function is:
\begin{align}
\min_{\mathbf{x}} \; H_0(\mathbf{x}) = N - \sum_{(i,j) \in E} x_{ij}.
\end{align}

Step 2 --- Write each constraint in its original algebraic form. The degree constraint requires that every vertex is incident to exactly two selected edges:
\begin{align}
\sum_{j \in V, j \neq i} x_{ij} = 2, \quad \forall i \in V.
\end{align}
(Subtour elimination constraints $\sum_{i \in S, j \in S, i < j} x_{ij} \le |S| - 1$ for all proper subsets $S \subset V$ are omitted from the penalty formulation here, as they are typically enforced via iterative cutting planes or are vacuous for small instances.)

Step 3 --- Convert each constraint to a penalty term. For a fixed vertex $i$, the squared penalty with weight $\lambda > 0$ is:
\begin{align}
P_i(\mathbf{x}) = \lambda \left( \sum_{j \in V, j \neq i} x_{ij} - 2 \right)^2.
\end{align}
Expanding the square and applying the idempotent property $x_{ij}^2 = x_{ij}$ for binary variables gives:
\begin{align}
P_i(\mathbf{x}) &= \lambda \left[ \left(\sum_{j \neq i} x_{ij}\right)^2 - 4 \sum_{j \neq i} x_{ij} + 4 \right] \\
&= \lambda \left[ \sum_{j \neq i} x_{ij} + \sum_{j \neq i} \sum_{\substack{k \neq i \\ k > j}} 2 x_{ij} x_{ik} - 4 \sum_{j \neq i} x_{ij} + 4 \right] \\
&= \lambda \left[ 4 - 3 \sum_{j \neq i} x_{ij} + 2 \sum_{j < k, \, j,k \neq i} x_{ij} x_{ik} \right].
\end{align}

Step 4 --- Combine objective and all penalty terms into the single QUBO Hamiltonian $H(\mathbf{x})$:
\begin{align}
H(\mathbf{x}) = N - \sum_{(i,j) \in E} x_{ij} + \sum_{i \in V} \lambda \left( \sum_{j \neq i} x_{ij} - 2 \right)^2.
\end{align}

Step 5 --- Read off the Q matrix entries and offset $c$ from $H$ in general closed form. Each variable $x_{ij}$ appears in exactly two degree constraints (for vertices $i$ and $j$). Collecting linear, quadratic, and constant terms yields:
\begin{align}
Q_{ij,ij} &= -1 - 6\lambda, \quad \forall (i,j) \in E, \\
Q_{ij,ik} &= 2\lambda, \quad \forall i \in V, \, j < k, \, (i,j), (i,k) \in E, \\
c &= N + 4N\lambda.
\end{align}
All other $Q$ entries are zero.

\paragraph{Penalty Weight Justification.}
The penalty weight $\lambda$ must be chosen sufficiently large to ensure that any violation of the degree constraints increases the total energy more than the maximum possible improvement in the objective function. The objective coefficient magnitude is $1$, and a single degree constraint can be violated by at most $2$ units (since degrees are sums of binary variables). To guarantee that feasible solutions are strictly preferred over infeasible ones, we require $\lambda > \max_{(i,j) \in E} |c_{ij}| \times \Delta_{\text{max}}$, where $\Delta_{\text{max}} = 2$ is the maximum constraint violation. Selecting $\lambda = 10$ satisfies this condition comfortably, ensuring the penalty dominates the objective landscape.

\paragraph{Correctness Argument.}
Minimizing $H(\mathbf{x})$ over $\{0, 1\}^{|E|}$ recovers the optimal Hamiltonian cycle because any feasible cycle satisfies all degree constraints exactly, yielding zero penalty energy. Infeasible assignments violate at least one degree constraint, incurring a penalty of at least $\lambda$. Since $\lambda$ is chosen to strictly dominate the objective coefficient range, any infeasible assignment has strictly higher energy than the best feasible assignment. Consequently, the global minimum of $H(\mathbf{x})$ must correspond to a degree-2 subgraph. (When combined with subtour elimination cuts or appropriate graph structures, this guarantees a single Hamiltonian cycle.)

\paragraph{Illustrative Example.}
Consider the complete graph $K_3$ with $N=3$ vertices $\{0, 1, 2\}$ and edges $E = \{(0,1), (0,2), (1,2)\}$. The decision variables are $x_{01}, x_{02}, x_{12}$. Substituting $N=3$ and $\lambda=10$ into the general formulas yields:
\begin{align}
H(\mathbf{x}) = 3 - (x_{01} + x_{02} + x_{12}) + 10 \sum_{i=0}^{2} \left( \sum_{j \neq i} x_{ij} - 2 \right)^2.
\end{align}
Expanding the penalty terms for each vertex gives:
\begin{align}
P_0 &= 40 - 30x_{01} - 30x_{02} + 20x_{01}x_{02}, \\
P_1 &= 40 - 30x_{01} - 30x_{12} + 20x_{01}x_{12}, \\
P_2 &= 40 - 30x_{02} - 30x_{12} + 20x_{02}x_{12}.
\end{align}
Summing the objective and penalties results in the concrete QUBO:
\begin{align}
H(\mathbf{x}) = 123 - 61(x_{01} + x_{02} + x_{12}) + 20(x_{01}x_{02} + x_{01}x_{12} + x_{02}x_{12}).
\end{align}
The corresponding Q matrix entries and offset are:
\begin{align}
Q_{00} = Q_{11} = Q_{22} = -61, \quad Q_{01} = Q_{02} = Q_{12} = 20, \quad c = 123.
\end{align}
The optimal assignment is $x_{01} = x_{02} = x_{12} = 1$, which selects all three edges and forms a valid triangle (Hamiltonian cycle). The minimum energy is $H(1,1,1) = 123 - 61(3) + 20(3) = 0$, confirming a zero missing-edge count.

\end{document}